\documentclass[11pt]{article}

\usepackage[letterpaper,margin=1in]{geometry}
\usepackage{microtype}
\usepackage{booktabs}
\usepackage{array}
\usepackage[table,dvipsnames]{xcolor}
\usepackage{enumitem}
\usepackage[most]{tcolorbox}
\usepackage[hidelinks]{hyperref}

\emergencystretch=3em  % absorb the occasional unbreakable prose line
\setlist{topsep=.45em,itemsep=.2em,parsep=0pt}
\makeatletter
\renewcommand\section{\@startsection{section}{1}{\z@}%
  {-2.2ex plus -.7ex minus -.2ex}{.9ex}%
  {\normalfont\Large\bfseries}}
\makeatother

% --- palette (matches the technical memo) ---
\definecolor{infoblue}{HTML}{2C3E50}
\definecolor{rule}{HTML}{BDC3C7}

% Plain-language callout, same visual language as the white papers.
\newtcolorbox{plain}{
  enhanced, breakable,
  colback=Sepia!7, colframe=Sepia!55, boxrule=0.4pt, arc=2pt,
  left=7pt, right=7pt, top=4pt, bottom=4pt,
  fontupper=\small,
  coltitle=Sepia!72!black, fonttitle=\bfseries\sffamily\footnotesize,
  title={THE WHOLE PIPELINE}
}

\title{\bfseries Your Scanner Says ``Critical.''\\ FedRAMP Asks a Harder Question.}
\author{Matthew Venne \\ \small Chief Technology Officer, stackArmor}
\date{July 2026}

\begin{document}
\maketitle

\begin{center}
\itshape A plain-language companion to ``A Deterministic, CVSS--Environmental
Method for FedRAMP VDR/VER Vulnerability Prioritization.''
\end{center}

\bigskip

\section*{The problem in one sentence}
A vulnerability scanner tells you how bad a flaw is \emph{in the abstract};
FedRAMP's new rules ask how bad it is \emph{for the agency whose data sits on that
specific box} --- and those are not the same question.

A ``Critical'' CVE on a throwaway sandbox can matter less than a ``Low'' CVE on the
database holding six agencies' records. Everyone in security knows this
intuitively. The hard part is turning that intuition into a number you can defend
in an audit, reproduce on every run, and explain to a 3PAO without hand-waving.

That's the gap this work fills.

\section*{What FedRAMP gave us, and what it left out}
FedRAMP's Vulnerability Detection and Response and Vulnerability Evaluation
Requirements (VDR/VER) introduce a clean prioritization vocabulary:
\begin{itemize}[leftmargin=1.5em]
\item \textbf{PAIN} --- \emph{Potential Agency Impact}, rated N1 through N5. This
  is the ``how bad for the agency'' score.
\item A \textbf{remediation matrix} --- how many days (or hours) you have to fix
  something, based on its PAIN level, how exploitable it is, and whether it's
  reachable from the internet.
\end{itemize}

What FedRAMP deliberately \emph{didn't} give us is the missing middle: the function
that takes a specific CVE on a specific asset and spits out an N-level. They
specified the output buckets and the deadlines, but not the classifier that fills
the buckets. Every provider is left to invent that themselves.

An invented, analyst-by-analyst classifier is the worst of both worlds: it's not
reproducible (two analysts disagree), and it's not defensible (you can't show your
work). So the question becomes: \textbf{can we build that classifier in a
principled, standardized way instead of making one up?}

\section*{The key insight: the classifier already exists}
Here's the move at the heart of the proposal. \textbf{The thing FedRAMP asks for
--- ``re-weight a vulnerability's impact by what's at stake on this asset'' --- is
exactly what the CVSS Environmental metrics were built to do.}

CVSS (the scoring system behind those ``Critical/High/Medium'' labels) has a
lesser-known \emph{Environmental} group designed for precisely this: it lets the
organization consuming a vulnerability re-score it based on how much \emph{their}
environment cares about confidentiality, integrity, and availability on the
affected system. Those knobs are called CR, IR, and AR --- the Confidentiality,
Integrity, and Availability \emph{Requirements}.

So the mapping is almost one-to-one:

\begin{center}
\begin{tabular}{@{}p{2.9in}p{2.9in}@{}}
\toprule
\textbf{What FedRAMP asks} & \textbf{What already answers it} \\
\midrule
How sensitive is this asset's data? & CVSS Security Requirements (CR / IR / AR) \\
How big is the potential agency impact? & CVSS Modified Impact Sub-Score \\
Is it likely to be exploited? & EPSS probability + KEV (known exploited) \\
Is it reachable from the internet? & Deployment/network context \\
One agency or many? & A scope multiplier on the result \\
\bottomrule
\end{tabular}
\end{center}

The payoff: \textbf{you don't need to invent a risk algebra.} You derive what
federal consequences are actually in scope, assign each asset its architectural
role, evaluate two facts about reachability and exploitation, and the published
CVSS math does the rest. The output is deterministic --- same inputs, same answer,
every analyst, every run.

\section*{How it actually works (the no-math version)}
Every vulnerability can hurt three things: \textbf{secrecy} (confidentiality),
\textbf{correctness} (integrity), and \textbf{uptime} (availability). For each one,
you multiply two numbers:
\begin{enumerate}[leftmargin=1.8em]
\item \emph{How badly does this flaw break it?} (from the CVE's own CVSS vector)
\item \emph{How much does this particular system depend on it?} (from the asset's
  security-impact profile)
\end{enumerate}

Combine those three into a single damage score from 0 to 1, then translate that
score into a plain word: \textbf{Minimal, Narrow, Disruptive, or Debilitating.} Add
one more fact --- does this asset serve one agency or many? --- and you land on an
N-level from N1 to N5.

That's the whole engine. The \emph{only} place human judgment enters the core
scoring is three cut-points that decide where one word turns into the next --- and
those are meant to be documented, governed, and back-tested against how your senior
analysts actually triage.

\section*{Before the asset score: what federal impact is actually possible?}
There is one important step before the arithmetic. A cloud service may be capable
of serving many agencies and many information types, but a particular deployment
has a narrower intended use. NIST SP~800-60 gives provisional confidentiality,
integrity, and availability profiles for federal information types. The agency
confirms which types it will actually use in the offering and whether any special
factors apply.

That produces a three-dimensional \textbf{Security Requirements ceiling}. Keep the
dimensions. If the system category is High because confidentiality is High while
integrity is Moderate and availability is Low, the ceiling is H/M/L --- not H/H/H.
A generic text box or file upload does not authorize every hypothetical kind of
federal data.

This is also where the point of view matters. PAIN asks about impact to the agency,
government operations, federal assets, and affected individuals. The CSP may care
more about the same event for brand or corporate-reputation reasons, but that
separate business risk does not silently raise the agency's PAIN inputs.

\section*{But where does the asset profile come from?}
Everything above hinges on that second set of numbers --- \emph{how much does this
system depend on secrecy, correctness, and uptime?} That's the uncapped CR/IR/AR
\textbf{asset security-impact profile}. The profile is the required input; the
method does not require one taxonomy for producing it.

A provider can assign the three objectives directly from an evidenced component
assessment, derive them through a role-based archetype catalog, or use another
governed mapping. The important rules are that CR, IR, and AR can differ, the same
evidence produces the same vector, and the record explains why each objective was
chosen.

\textbf{Asset archetypes remain a strong example.} Instead of hand-rating every
server in isolation, a provider can sort assets into a short catalog of roles based
on \textbf{what they do inside the service} --- a database, a login service, a
cache, a message bus, a CI/CD runner, a public load balancer, and so on. Each role
maps to a pre-filled and justified profile. Tag an asset with its archetype, and the
profile is inherited automatically. The original PAIN paper uses this approach in
its worked architectures because it makes the examples concrete; it is not a
required CSP taxonomy.

The benefit is the \textbf{auditable derivation trail}. A bare ``CR: Low, IR: High,
AR: Low'' assignment can be hard to reconstruct six months later. A
\emph{data-backbone message broker} role makes one possible rationale visible and
consistent with comparable brokers. A direct assignment can be just as defensible
when it retains the architecture, trust, privilege, dependency, data-handling, and
consequence evidence behind each objective.

For providers using archetypes, classify by role, not software kind. The same
in-memory store can be a throwaway cache, a session-token store, or a job broker;
\emph{how it is used} drives the profile, not \emph{what product it is}. The example
catalog recommends a ``control-plane lens first'' rule: if an asset can deploy,
orchestrate, hold cross-estate credentials, or actuate configuration, classify it
by that powerful control function regardless of the data it happens to store.

A scalar asset-value label is a permitted shortcut: Low becomes L/L/L, Moderate
becomes M/M/M, and High becomes H/H/H. It is also deliberately cautioned against
as the only interface. It forces all three objectives to move together and records
the conclusion with less evidence, so an implementation should always permit an
independent CR/IR/AR profile.

So the full pipeline reads top to bottom like this:

\begin{plain}
\textbf{Derive intended-use ceiling} --- which federal information types and
consequences are in scope $\rightarrow$ \textbf{resolve the asset security-impact
profile} --- by direct assessment, an archetype, or another governed method
$\rightarrow$ \textbf{take the lower value
on each objective} --- effective CR/IR/AR for confidentiality, integrity, and
availability $\rightarrow$ \textbf{feed the CVSS
Environmental impact math} --- re-weighting the CVE's raw impact by what's at stake
\emph{here} $\rightarrow$ \textbf{produces the PAIN level (N1--N5)} --- plus, with
the reachability and exploitation facts, a concrete deadline.
\end{plain}

Two important guardrails on any profile derivation:
\begin{itemize}[leftmargin=1.5em]
\item \textbf{The archetype list is an example, not a standard.} The proposal
  publishes a sample catalog as an \emph{existence proof}. Every provider may use
  its own taxonomy, direct objective assessment, or another governed mapping.
  Agency FIPS-199 and information-type context belongs in the separate ceiling.
\item \textbf{Mislabeling is a downgrade attack.} Because the whole score rides on
  the profile, assigning an asset \emph{down} makes its flaws look smaller than
  they are. Profile assignments, mappings, and overrides should be controlled and
  reviewed like any other security setting --- and unresolved assets fail loud at
  H/H/H (see below), never quiet.
\end{itemize}

\section*{The same flaw, two very different verdicts}
This is the part worth internalizing, because it's the entire point of the method.
Take one identical remote-code-execution CVE:
\begin{itemize}[leftmargin=1.5em]
\item \textbf{On a crown-jewel datastore whose security-impact profile and
  intended-use ceiling retain the necessary High requirements, holding multiple
  agencies' data:} scores
  Debilitating, multi-agency $\rightarrow$ \textbf{N5}, remediate in hours.
\item \textbf{On a single-agency throwaway sandbox:} same CVE, scores Disruptive
  $\rightarrow$ \textbf{N3}, a much longer clock.
\end{itemize}

The vulnerability didn't change. The \emph{asset} did all the work. And a
``Low''-severity availability flaw on a high-uptime message backbone can
legitimately land at \textbf{N4 or N5} --- decoupled entirely from the vendor's
qualitative label.

\section*{The PAIN level says how \emph{bad}. Two more facts say how \emph{fast}.}
The N-level tells you the potential impact, but not your deadline. Two real-world
facts tighten the remediation clock:
\begin{itemize}[leftmargin=1.5em]
\item \textbf{Is it actually being exploited --- or likely to be?} This is the
  \emph{exploitability} axis. A flaw is treated as likely-exploitable if it clears
  an EPSS probability threshold (EPSS being the statistical ``how likely is this to
  be exploited in the wild'' score) \textbf{or} if it's known to be actively
  exploited --- for example, listed in CISA's KEV catalog. Either signal is enough.
  Whichever predicate fires moves the finding to a faster remediation column.
\item \textbf{Is it reachable from the internet?} A vulnerable surface an
  unauthenticated outside attacker can actually reach --- through a public IP, a
  public load balancer, an ingress, or any edge path that doesn't authenticate
  first --- is far more urgent than the same flaw buried behind authentication on
  an internal-only service. (A route gated by an identity-aware proxy in a
  \emph{remote-access} role --- your own people, with the application's own login
  still behind it --- counts as \emph{not} internet-reachable: the attacker can't
  get to the vulnerable code.)
\end{itemize}

Stack those up and the logic is intuitive: \textbf{more severe + more exploitable +
more exposed = less time to fix.} A lookup table turns the combination into a
concrete deadline, with tighter clocks for higher-assurance providers. So an
internet-facing, actively-exploited flaw on a shared multi-agency edge might be a
12-hour fix, while the same-tier flaw that's sub-threshold on EPSS and tucked
behind authentication gets a much longer runway.

There's also an optional refinement that can raise the \emph{level} rather than the
clock: where an authoritative source reports that exploiting a flaw grants
\emph{total} technical impact, the affected dimensions can be floored to High
before scoring --- pushing a weak-looking CVE up a tier. It never invents impact on
a dimension the CVE doesn't touch.

\section*{Fail loud, not quiet}
A critical design choice: \textbf{missing information never makes a problem look
smaller.} An asset whose security-impact profile is unresolved defaults to H/H/H
--- so it scores loudly and surfaces itself for assessment, rather than hiding in
the noise. ``Unknown'' is treated as ``serious until proven otherwise.''

\section*{Scoring isn't the end --- disposition is}
A scanner hit isn't a confirmed problem. The flagged code might not be present,
might never run, might need a configuration nobody set, or might already be
neutralized by a control. Sorting that out is the \textbf{disposition} step, and it
can only happen \emph{after} a finding is detected.

The method records dispositions using \textbf{VEX} (Vulnerability Exploitability
eXchange) --- a standardized, signable, machine-readable artifact. Crucially,
disposition is an \emph{overlay, not a re-score}: a finding's computed PAIN is
always retained on record even when it's suppressed. That makes every ``this
doesn't apply to us'' decision \textbf{auditable, attributable, and reversible} ---
with stronger evidence required to wave off a high-stakes finding than a trivial
internal one.

\section*{Why this matters for the bigger picture}
FedRAMP's VDR/VER rule is its implementation of CISA's Binding Operational
Directive 26-04 --- the broader push to move remediation prioritization \emph{off}
raw CVSS base scores and \emph{onto} exploitability and mission impact. FedRAMP
kept the directive's philosophy but swapped in its own instruments (PAIN in place
of SSVC, among others). Because it chose PAIN, an off-the-shelf SSVC tool doesn't
satisfy the rule by itself --- and the classifier that produces an N-level was
exactly the missing piece. This method supplies it in standards-grounded terms.

\section*{The bottom line}
You can summarize the entire proposal in four claims:
\begin{enumerate}[leftmargin=1.8em]
\item \textbf{FedRAMP specified the vocabulary and the deadlines but not the
  classifier.} Everyone has to build one.
\item \textbf{You don't have to invent it.} CVSS Environmental metrics already
  encode ``re-weight impact by what's at stake'' --- which \emph{is} the
  agency-impact question.
\item \textbf{With governed inputs} --- an intended-use ceiling, the asset's
  security-impact profile, and a one-agency-or-many affected scope --- every finding gets a
  deterministic, auditable N-level and a concrete deadline. The inputs change
  when use, architecture, or tenancy changes; the finding math does not.
\item \textbf{It's defensible.} Same inputs, same output, every time --- with the
  full derivation attached to each finding, and a fail-safe that errs toward
  over-classifying rather than under.
\end{enumerate}

The deeper paper has the worked examples, the exact formulas, the reference
architectures, and a sample archetype-to-profile catalog. The catalog is one
illustrative derivation method, \emph{not} a required taxonomy or standard. But the
idea you came for
is simple: \textbf{stop asking ``how bad is this CVE?'' and start asking ``how bad
is this CVE \emph{here}?'' --- and let a formula everyone can check answer it the
same way every time.}

\bigskip
\noindent\textcolor{rule}{\rule{\textwidth}{0.4pt}}
\medskip

\noindent\emph{Read the full technical memo:
\href{https://stackarmor.github.io/rfc-fedramp-vdr/vdr-pain-cvss.pdf}{A
Deterministic, CVSS--Environmental Method for FedRAMP VDR/VER Vulnerability
Prioritization}. Read the upstream derivation:
\href{https://stackarmor.github.io/rfc-fedramp-vdr/vdr-pain-security-requirements.pdf}{Before
the PAIN Equation: Deriving Security-Requirements Ceilings from Intended Federal
Information Types}.}

\end{document}
